\subsection{扩张的超越次数}

定理3. --- 设 E 是域 K 的一个扩域。 E 在 K 上的所有超越基都具有相同的基数。

只需证明不等式 \(\operatorname{Card}(\mathbf{B}) \geqslant \operatorname{Card}(\mathbf{B}')\) ，其中 \(\mathbf{B}\) 与 \(\mathbf{B}'\) 是 \(\mathbf{E}\) 在 \(\mathbf{K}\) 上的两个超越基，并且我们可设 \(\mathbf{B}'\) 非空。我们先假设 \(\mathbf{B}\) 有限，并对其基数 \(n\) 作归纳；当 \(n = 0\) 时， \(\mathbf{E}\) 在 \(\mathbf{K}\) 上是代数的，且 \(\mathbf{B}'\) 为空。现设 \(n \geqslant 1\) ；给定 \(x \in \mathbf{B}'\) ，交换定理提供子集 \(\mathbf{C}\) （ \(\mathbf{B}\) 的），使得 \(x \notin \mathbf{C}\) 且 \(\{x\} \cup \mathbf{C}\) 是 \(\mathbf{E}\) 在 \(\mathbf{K}\) 上的一个超越基；由命题7我们有 \(\mathbf{C} \neq \mathbf{B}\) ，因此 \(\operatorname{Card} \mathbf{C} < n\) 。令 \(\mathbf{K}_1 = \mathbf{K}(x)\) 且 \(\mathbf{C}' = \mathbf{B}' - \{x\}\) ；则 \(\mathbf{C}\) 与 \(\mathbf{C}'\) 在域 \(\mathbf{K}_1\) 上代数自由（V，第107页，命题4），且 \(\mathbf{E}\) 既在 \(\mathbf{K}_1(\mathbf{C}) = \mathbf{K}(\mathbf{C} \cup \{x\})\) 上、又在 \(\mathbf{K}_1(\mathbf{C}') = \mathbf{K}(\mathbf{B}')\) 上是代数的。换言之， \(\mathbf{C}\) 与 \(\mathbf{C}'\) 是 \(\mathbf{E}\) 在 \(\mathbf{K}_1\) 上的两个超越基。由于 \(\operatorname{Card}(\mathbf{C}) < n\) ，归纳假设给出不等式 \(\operatorname{Card}(\mathbf{C}') \leqslant \operatorname{Card}(\mathbf{C}) \leqslant n - 1\) ，从而 \(\operatorname{Card}(\mathbf{B}') \leqslant n = \operatorname{Card}(\mathbf{B})\) 。

现设 B 是无限的。每个 \(x \in B\) 在 \(\mathrm{K}(\mathrm{B}')\) 上是代数的，于是存在有限子集 \(\mathrm{S}(x)\) （ \(B'\) 的），使得 x 在 \(\mathrm{K}(\mathrm{S}(x))\) 上是代数的。记 \(\mathrm{S} = \bigcup_{x \in \mathrm{B}} \mathrm{S}(x)\) ，则 \(S \subset B'\) ，且由于 B 是无限的，我们有 \(\operatorname{Card}(\mathrm{S}) \leqslant \operatorname{Card}(\mathrm{B})\)

（E，III，第49页，推论3）。但由于 B 的每个元素在 K(S) 上是代数的，且 E 在 K(B) 上是代数的，我们得出 E 在 K(S) 上是代数的（V，第19页，命题3）。现在命题7蕴含 S = B'，从而得到所希望的不等式 Card(B') ≤ Card(B)。

定义4. --- 设 E 是域 K 的一个扩域。 E 在 K 上的每个超越基的基数称为 E 在 K 上的超越次数，记作 tr. deg \(_{K}\) E。

定理2、定理3以及定义4蕴含以下推论，其中 E 表示域 K 的一个具有有限超越次数 n 的扩域。

推论1. --- 设 S 是 E 的一个子集，使得 E 在 K(S) 上是代数的。则 Card(S) ≥ n；若 S 的基数等于 n，则 S 在 K 上是代数自由的（从而此时它是 E 在 K 上的一个超越基）。

推论2. --- 假设 \(\mathrm{E} = \mathrm{K}(x_1, \ldots, x_m)\) ，则 \(m \geqslant n\) ；此外若 \(m = n\) ，则 \((x_1, \ldots, x_m)\) 是 \(\mathrm{E}\) 在 \(\mathrm{K}\) 上的一个纯基，从而 \(\mathrm{E}\) 是 \(\mathrm{K}\) 的一个纯扩张。

推论3. --- E 的每个在 K 上代数自由的子集至多有 n 个元素，若它恰好有 n 个元素，则它是 E 在 K 上的一个超越基。

定理4. --- 设 K、E 和 F 是三个域，满足 \(K \subset E \subset F\) 。若 S 是 E 在 K 上的一个超越基，T 是 F 在 E 上的一个超越基，则 \(S \cap T\) 为空，且 \(S \cup T\) 是 F 在 K 上的一个超越基。

因为 F 在 \(E(T)\) 上是代数的；此外， \(E(T)\) 在域 \(K(S \cup T) = K(S)(T)\) 上是代数的，因为 E 在 \(K(S)\) 上是代数的（V，第18页，推论2）；因此（V，第19页，命题3）F 在 \(K(S \cup T)\) 上是代数的。另一方面，T 在 E 上代数自由，从而在 \(K(S)\) 上更是如此，故（V，第107页，命题4） \(S \cup T\) 在 K 上代数自由，且 \(S \cap T = \varnothing\) 。

推论. --- 设 K、E 和 F 是三个域，满足 \(K \subset E \subset F\) 。则我们有

\[
\operatorname{tr} \cdot \deg_ {K} F = \operatorname{tr} \cdot \deg_ {K} E + \operatorname{tr} \cdot \deg_ {E} F.\tag{1}
\]

